% Options for packages loaded elsewhere
\PassOptionsToPackage{unicode}{hyperref}
\PassOptionsToPackage{hyphens}{url}
%
\documentclass[
]{article}
\usepackage{amsmath,amssymb}
\usepackage{iftex}
\ifPDFTeX
  \usepackage[T1]{fontenc}
  \usepackage[utf8]{inputenc}
  \usepackage{textcomp} % provide euro and other symbols
\else % if luatex or xetex
  \usepackage{unicode-math} % this also loads fontspec
  \defaultfontfeatures{Scale=MatchLowercase}
  \defaultfontfeatures[\rmfamily]{Ligatures=TeX,Scale=1}
\fi
\usepackage{lmodern}
\ifPDFTeX\else
  % xetex/luatex font selection
\fi
% Use upquote if available, for straight quotes in verbatim environments
\IfFileExists{upquote.sty}{\usepackage{upquote}}{}
\IfFileExists{microtype.sty}{% use microtype if available
  \usepackage[]{microtype}
  \UseMicrotypeSet[protrusion]{basicmath} % disable protrusion for tt fonts
}{}
\makeatletter
\@ifundefined{KOMAClassName}{% if non-KOMA class
  \IfFileExists{parskip.sty}{%
    \usepackage{parskip}
  }{% else
    \setlength{\parindent}{0pt}
    \setlength{\parskip}{6pt plus 2pt minus 1pt}}
}{% if KOMA class
  \KOMAoptions{parskip=half}}
\makeatother
\usepackage{xcolor}
\usepackage[margin=1in]{geometry}
\usepackage{color}
\usepackage{fancyvrb}
\newcommand{\VerbBar}{|}
\newcommand{\VERB}{\Verb[commandchars=\\\{\}]}
\DefineVerbatimEnvironment{Highlighting}{Verbatim}{commandchars=\\\{\}}
% Add ',fontsize=\small' for more characters per line
\usepackage{framed}
\definecolor{shadecolor}{RGB}{248,248,248}
\newenvironment{Shaded}{\begin{snugshade}}{\end{snugshade}}
\newcommand{\AlertTok}[1]{\textcolor[rgb]{0.94,0.16,0.16}{#1}}
\newcommand{\AnnotationTok}[1]{\textcolor[rgb]{0.56,0.35,0.01}{\textbf{\textit{#1}}}}
\newcommand{\AttributeTok}[1]{\textcolor[rgb]{0.13,0.29,0.53}{#1}}
\newcommand{\BaseNTok}[1]{\textcolor[rgb]{0.00,0.00,0.81}{#1}}
\newcommand{\BuiltInTok}[1]{#1}
\newcommand{\CharTok}[1]{\textcolor[rgb]{0.31,0.60,0.02}{#1}}
\newcommand{\CommentTok}[1]{\textcolor[rgb]{0.56,0.35,0.01}{\textit{#1}}}
\newcommand{\CommentVarTok}[1]{\textcolor[rgb]{0.56,0.35,0.01}{\textbf{\textit{#1}}}}
\newcommand{\ConstantTok}[1]{\textcolor[rgb]{0.56,0.35,0.01}{#1}}
\newcommand{\ControlFlowTok}[1]{\textcolor[rgb]{0.13,0.29,0.53}{\textbf{#1}}}
\newcommand{\DataTypeTok}[1]{\textcolor[rgb]{0.13,0.29,0.53}{#1}}
\newcommand{\DecValTok}[1]{\textcolor[rgb]{0.00,0.00,0.81}{#1}}
\newcommand{\DocumentationTok}[1]{\textcolor[rgb]{0.56,0.35,0.01}{\textbf{\textit{#1}}}}
\newcommand{\ErrorTok}[1]{\textcolor[rgb]{0.64,0.00,0.00}{\textbf{#1}}}
\newcommand{\ExtensionTok}[1]{#1}
\newcommand{\FloatTok}[1]{\textcolor[rgb]{0.00,0.00,0.81}{#1}}
\newcommand{\FunctionTok}[1]{\textcolor[rgb]{0.13,0.29,0.53}{\textbf{#1}}}
\newcommand{\ImportTok}[1]{#1}
\newcommand{\InformationTok}[1]{\textcolor[rgb]{0.56,0.35,0.01}{\textbf{\textit{#1}}}}
\newcommand{\KeywordTok}[1]{\textcolor[rgb]{0.13,0.29,0.53}{\textbf{#1}}}
\newcommand{\NormalTok}[1]{#1}
\newcommand{\OperatorTok}[1]{\textcolor[rgb]{0.81,0.36,0.00}{\textbf{#1}}}
\newcommand{\OtherTok}[1]{\textcolor[rgb]{0.56,0.35,0.01}{#1}}
\newcommand{\PreprocessorTok}[1]{\textcolor[rgb]{0.56,0.35,0.01}{\textit{#1}}}
\newcommand{\RegionMarkerTok}[1]{#1}
\newcommand{\SpecialCharTok}[1]{\textcolor[rgb]{0.81,0.36,0.00}{\textbf{#1}}}
\newcommand{\SpecialStringTok}[1]{\textcolor[rgb]{0.31,0.60,0.02}{#1}}
\newcommand{\StringTok}[1]{\textcolor[rgb]{0.31,0.60,0.02}{#1}}
\newcommand{\VariableTok}[1]{\textcolor[rgb]{0.00,0.00,0.00}{#1}}
\newcommand{\VerbatimStringTok}[1]{\textcolor[rgb]{0.31,0.60,0.02}{#1}}
\newcommand{\WarningTok}[1]{\textcolor[rgb]{0.56,0.35,0.01}{\textbf{\textit{#1}}}}
\usepackage{graphicx}
\makeatletter
\newsavebox\pandoc@box
\newcommand*\pandocbounded[1]{% scales image to fit in text height/width
  \sbox\pandoc@box{#1}%
  \Gscale@div\@tempa{\textheight}{\dimexpr\ht\pandoc@box+\dp\pandoc@box\relax}%
  \Gscale@div\@tempb{\linewidth}{\wd\pandoc@box}%
  \ifdim\@tempb\p@<\@tempa\p@\let\@tempa\@tempb\fi% select the smaller of both
  \ifdim\@tempa\p@<\p@\scalebox{\@tempa}{\usebox\pandoc@box}%
  \else\usebox{\pandoc@box}%
  \fi%
}
% Set default figure placement to htbp
\def\fps@figure{htbp}
\makeatother
\setlength{\emergencystretch}{3em} % prevent overfull lines
\providecommand{\tightlist}{%
  \setlength{\itemsep}{0pt}\setlength{\parskip}{0pt}}
\setcounter{secnumdepth}{-\maxdimen} % remove section numbering
\usepackage{bookmark}
\IfFileExists{xurl.sty}{\usepackage{xurl}}{} % add URL line breaks if available
\urlstyle{same}
\hypersetup{
  pdftitle={Time Series CV CO2},
  hidelinks,
  pdfcreator={LaTeX via pandoc}}

\title{Time Series CV CO2}
\author{}
\date{\vspace{-2.5em}2026-06-24}

\begin{document}
\maketitle

\begin{Shaded}
\begin{Highlighting}[]
\FunctionTok{library}\NormalTok{(forecast)}
\end{Highlighting}
\end{Shaded}

\begin{verbatim}
## Warning: package 'forecast' was built under R version 4.5.3
\end{verbatim}

\begin{Shaded}
\begin{Highlighting}[]
\NormalTok{x }\OtherTok{=}\NormalTok{ co2}
\CommentTok{\# split the last one year for testing }
\NormalTok{test}\OtherTok{=} \FunctionTok{window}\NormalTok{(x,}\AttributeTok{start=}\FunctionTok{c}\NormalTok{(}\DecValTok{1997}\NormalTok{,}\DecValTok{1}\NormalTok{))}
\NormalTok{train }\OtherTok{=} \FunctionTok{window}\NormalTok{(x,}\AttributeTok{end =} \FunctionTok{c}\NormalTok{(}\DecValTok{1996}\NormalTok{,}\DecValTok{12}\NormalTok{))}
\FunctionTok{autoplot}\NormalTok{(train, }\AttributeTok{series=}\StringTok{"Train"}\NormalTok{)}\SpecialCharTok{+}
  \FunctionTok{autolayer}\NormalTok{(test,}\AttributeTok{series =} \StringTok{"Test"}\NormalTok{)}
\end{Highlighting}
\end{Shaded}

\pandocbounded{\includegraphics[keepaspectratio]{ts_co2_cv_files/figure-latex/unnamed-chunk-1-1.pdf}}

\section{Cross validation to compare different
model}\label{cross-validation-to-compare-different-model}

\begin{Shaded}
\begin{Highlighting}[]
\NormalTok{mape }\OtherTok{=} \ControlFlowTok{function}\NormalTok{(forecast,truth)\{}
  \FunctionTok{return}\NormalTok{ (}\DecValTok{100}\SpecialCharTok{*}\FunctionTok{mean}\NormalTok{(}\FunctionTok{abs}\NormalTok{(truth}\SpecialCharTok{{-}}\NormalTok{forecast)}\SpecialCharTok{/}\NormalTok{truth))}
\NormalTok{\}}
\CommentTok{\#initialize datafram to store cv result of models }
\NormalTok{cv\_results }\OtherTok{=} \FunctionTok{data.frame}\NormalTok{(}
  \AttributeTok{Model =} \FunctionTok{character}\NormalTok{(),}
  \AttributeTok{Fold =} \FunctionTok{integer}\NormalTok{(),}
  \AttributeTok{MAPE =} \FunctionTok{numeric}\NormalTok{(),}
  \AttributeTok{stringsAsFactors =} \ConstantTok{FALSE}

\NormalTok{)}

\NormalTok{start\_year }\OtherTok{=} \FunctionTok{start}\NormalTok{(train)[}\DecValTok{1}\NormalTok{]}
\NormalTok{end\_year }\OtherTok{=} \FunctionTok{end}\NormalTok{(train)[}\DecValTok{1}\NormalTok{]}

\ControlFlowTok{for}\NormalTok{ (i }\ControlFlowTok{in}\NormalTok{ (start\_year}\SpecialCharTok{+}\DecValTok{1}\NormalTok{)}\SpecialCharTok{:}\NormalTok{(end\_year[}\DecValTok{1}\NormalTok{]}\SpecialCharTok{{-}}\DecValTok{1}\NormalTok{))\{}
  \CommentTok{\# the first train fold has 24 observations}
\NormalTok{  train\_fold }\OtherTok{=} \FunctionTok{window}\NormalTok{(train,}\AttributeTok{end =} \FunctionTok{c}\NormalTok{(i,}\DecValTok{12}\NormalTok{))}
  \CommentTok{\# the validation fold has 12 observation }
\NormalTok{  val\_fold }\OtherTok{=} \FunctionTok{window}\NormalTok{(train,}\AttributeTok{start=}\FunctionTok{c}\NormalTok{(i}\SpecialCharTok{+}\DecValTok{1}\NormalTok{,}\DecValTok{1}\NormalTok{),}\AttributeTok{end =} \FunctionTok{c}\NormalTok{(i}\SpecialCharTok{+}\DecValTok{1}\NormalTok{,}\DecValTok{12}\NormalTok{))}
  

  \CommentTok{\#Simple exponential smoothing}
\NormalTok{  ses }\OtherTok{=} \FunctionTok{HoltWinters}\NormalTok{(train\_fold,}\AttributeTok{alpha =} \ConstantTok{NULL}\NormalTok{,}\AttributeTok{beta =} \ConstantTok{FALSE}\NormalTok{, }\AttributeTok{gamma=}\ConstantTok{FALSE}\NormalTok{)}
\NormalTok{  ses\_pred }\OtherTok{=} \FunctionTok{predict}\NormalTok{(ses,}\AttributeTok{n.ahead=}\FunctionTok{length}\NormalTok{(val\_fold),)}
\NormalTok{  cv\_results }\OtherTok{=} \FunctionTok{rbind}\NormalTok{( }
\NormalTok{                cv\_results,}
                \FunctionTok{data.frame}\NormalTok{(}
                  \AttributeTok{Model =} \StringTok{"SES"}\NormalTok{,}
                  \AttributeTok{Fold =}\NormalTok{ i,}
                  \AttributeTok{MAPE =} \FunctionTok{mape}\NormalTok{(ses\_pred[,}\StringTok{"fit"}\NormalTok{],val\_fold)}
\NormalTok{                )}
\NormalTok{                )}
  \CommentTok{\# non seasonal exponential smooth}
\NormalTok{  non\_seasonal }\OtherTok{=} \FunctionTok{HoltWinters}\NormalTok{(train\_fold,}\AttributeTok{alpha =} \ConstantTok{NULL}\NormalTok{,}\AttributeTok{beta =} \ConstantTok{NULL}\NormalTok{,}\AttributeTok{gamma =} \ConstantTok{FALSE}\NormalTok{)}
\NormalTok{  non\_seasonal\_pred }\OtherTok{=} \FunctionTok{predict}\NormalTok{(non\_seasonal,}\AttributeTok{n.ahead=}\FunctionTok{length}\NormalTok{(val\_fold))}
\NormalTok{  cv\_results }\OtherTok{=} \FunctionTok{rbind}\NormalTok{( }
\NormalTok{                cv\_results,}
                \FunctionTok{data.frame}\NormalTok{(}
                  \AttributeTok{Model =} \StringTok{"non\_seasonal"}\NormalTok{,}
                  \AttributeTok{Fold =}\NormalTok{ i,}
                  \AttributeTok{MAPE =} \FunctionTok{mape}\NormalTok{(non\_seasonal\_pred[,}\StringTok{"fit"}\NormalTok{],val\_fold)}
\NormalTok{                )}
\NormalTok{                )}
  \CommentTok{\#additive seasonal exponential smoothing }
\NormalTok{  add\_seasonal }\OtherTok{=} \FunctionTok{HoltWinters}\NormalTok{(train\_fold,}\AttributeTok{alpha =} \ConstantTok{NULL}\NormalTok{,}\AttributeTok{beta =} \ConstantTok{NULL}\NormalTok{,}\AttributeTok{gamma =} \ConstantTok{NULL}\NormalTok{)}
\NormalTok{  add\_seasonal\_pred }\OtherTok{=} \FunctionTok{predict}\NormalTok{(add\_seasonal,}\AttributeTok{n.ahead=}\FunctionTok{length}\NormalTok{(val\_fold))}
\NormalTok{  cv\_results }\OtherTok{=} \FunctionTok{rbind}\NormalTok{( }
\NormalTok{                cv\_results,}
                \FunctionTok{data.frame}\NormalTok{(}
                  \AttributeTok{Model =} \StringTok{"add\_seasonal"}\NormalTok{,}
                  \AttributeTok{Fold =}\NormalTok{ i,}
                  \AttributeTok{MAPE =} \FunctionTok{mape}\NormalTok{(add\_seasonal\_pred[,}\StringTok{"fit"}\NormalTok{],val\_fold)}
\NormalTok{                )}
\NormalTok{                )}
  \CommentTok{\#multiplicative seasonal }
\NormalTok{  multi\_seasonal }\OtherTok{=} \FunctionTok{HoltWinters}\NormalTok{(train\_fold,}\AttributeTok{alpha =} \ConstantTok{NULL}\NormalTok{,}\AttributeTok{beta =} \ConstantTok{NULL}\NormalTok{,}\AttributeTok{gamma =} \ConstantTok{NULL}\NormalTok{, }\AttributeTok{seasonal =} \StringTok{"multi"}\NormalTok{)}
\NormalTok{  multi\_seasonal\_pred }\OtherTok{=} \FunctionTok{predict}\NormalTok{(multi\_seasonal,}\AttributeTok{n.ahead=}\FunctionTok{length}\NormalTok{(val\_fold))}
\NormalTok{  cv\_results }\OtherTok{=} \FunctionTok{rbind}\NormalTok{( }
\NormalTok{              cv\_results,}
              \FunctionTok{data.frame}\NormalTok{(}
                \AttributeTok{Model =} \StringTok{"multi\_seasonal"}\NormalTok{,}
                \AttributeTok{Fold =}\NormalTok{ i,}
                \AttributeTok{MAPE =} \FunctionTok{mape}\NormalTok{(multi\_seasonal\_pred[,}\StringTok{"fit"}\NormalTok{],val\_fold)}
\NormalTok{              )}
\NormalTok{              )}

\NormalTok{\}}
\end{Highlighting}
\end{Shaded}

\begin{verbatim}
## Warning in HoltWinters(train_fold, alpha = NULL, beta = NULL, gamma = FALSE):
## optimization difficulties: ERROR: ABNORMAL_TERMINATION_IN_LNSRCH
## Warning in HoltWinters(train_fold, alpha = NULL, beta = NULL, gamma = FALSE):
## optimization difficulties: ERROR: ABNORMAL_TERMINATION_IN_LNSRCH
\end{verbatim}

\begin{verbatim}
## Warning in HoltWinters(train_fold, alpha = NULL, beta = NULL, gamma = NULL, :
## optimization difficulties: ERROR: ABNORMAL_TERMINATION_IN_LNSRCH
\end{verbatim}

\begin{Shaded}
\begin{Highlighting}[]
\NormalTok{summary\_results }\OtherTok{\textless{}{-}} \FunctionTok{do.call}\NormalTok{(}
\NormalTok{  data.frame,}
  \FunctionTok{aggregate}\NormalTok{(}
\NormalTok{    MAPE }\SpecialCharTok{\textasciitilde{}}\NormalTok{ Model,}
    \AttributeTok{data =}\NormalTok{ cv\_results,}
    \AttributeTok{FUN =} \ControlFlowTok{function}\NormalTok{(x)}
      \FunctionTok{c}\NormalTok{(}
        \AttributeTok{Mean\_MAPE =} \FunctionTok{mean}\NormalTok{(x, }\AttributeTok{na.rm =} \ConstantTok{TRUE}\NormalTok{),}
        \AttributeTok{SD\_MAPE   =} \FunctionTok{sd}\NormalTok{(x, }\AttributeTok{na.rm =} \ConstantTok{TRUE}\NormalTok{)}
\NormalTok{      )}
\NormalTok{  )}
\NormalTok{)}

\FunctionTok{print}\NormalTok{(summary\_results)}
\end{Highlighting}
\end{Shaded}

\begin{verbatim}
##            Model MAPE.Mean_MAPE MAPE.SD_MAPE
## 1   add_seasonal      0.1135635   0.05517814
## 2 multi_seasonal      0.1109951   0.05256529
## 3   non_seasonal      1.3978538   0.72163382
## 4            SES      0.6144831   0.04994032
\end{verbatim}

\begin{Shaded}
\begin{Highlighting}[]
\CommentTok{\# train the best model with the full trianing set }
\NormalTok{best\_model }\OtherTok{=} \FunctionTok{HoltWinters}\NormalTok{(train,}\AttributeTok{alpha =} \ConstantTok{NULL}\NormalTok{,}\AttributeTok{beta =} \ConstantTok{NULL}\NormalTok{,}\AttributeTok{gamma =} \ConstantTok{NULL}\NormalTok{, }\AttributeTok{seasonal =} \StringTok{"multi"}\NormalTok{)}
\CommentTok{\# test on the left out dataset}
\NormalTok{best\_forecast }\OtherTok{=} \FunctionTok{predict}\NormalTok{(best\_model,}\AttributeTok{n.ahead=}\FunctionTok{length}\NormalTok{(test))}
\FunctionTok{mape}\NormalTok{(best\_forecast[,}\StringTok{"fit"}\NormalTok{],test)}
\end{Highlighting}
\end{Shaded}

\begin{verbatim}
## [1] 0.09558458
\end{verbatim}

\begin{Shaded}
\begin{Highlighting}[]
\FunctionTok{autoplot}\NormalTok{(train, }\AttributeTok{series=}\StringTok{"Train"}\NormalTok{)}\SpecialCharTok{+}
  \FunctionTok{autolayer}\NormalTok{(test,}\AttributeTok{series =} \StringTok{"Test"}\NormalTok{)}\SpecialCharTok{+}
  \FunctionTok{autolayer}\NormalTok{(best\_forecast,}\AttributeTok{series =} \StringTok{"Forecast"}\NormalTok{)}
\end{Highlighting}
\end{Shaded}

\pandocbounded{\includegraphics[keepaspectratio]{ts_co2_cv_files/figure-latex/unnamed-chunk-4-1.pdf}}

\end{document}
